\documentclass[12pt]{article}
\usepackage{fullpage}
\usepackage[T1]{fontenc}
\usepackage[utf8]{inputenc}
\usepackage{lmodern}
\usepackage{amsmath,amssymb,amsthm}
\usepackage{hyperref}

\newtheorem{theorem}{Theorem}
\newtheorem{lemma}{Lemma}
\newtheorem{proposition}{Proposition}
\newtheorem{corollary}{Corollary}
\theoremstyle{remark}
\newtheorem{remark}{Remark}

\newcommand{\leg}[2]{\left(\frac{#1}{#2}\right)}
\newcommand{\bin}[2]{\binom{#1}{#2}}

\begin{document}

\title{On the supercongruences for $\sum (3k+1)\binom{2k}{k}^3/16^k$:\\
established facts, a refuted statement, and reductions}
\author{}
\date{}
\maketitle

\begin{abstract}
We analyse the five assertions (i)--(v).  We prove in full an elementary
\emph{linking identity} (Theorem~\ref{thm:link}) which ties together the two
congruences in each of (i) and (ii); using it together with the classical
mod\,$p^2$ evaluation of $\sum\binom{2k}{k}^3/16^k$ we deduce part~(ii)
completely and the first congruence of part~(i).  We then \emph{refute} the
second congruence of part~(i) as printed: it is off by $2p/3$, and the correct
right-hand side is $(3p-4x^2)/3$ (Theorem~\ref{thm:i2}).  For the remaining deep
parts (iii), (iv), (v) we give precise statements and exhibit exactly how they
reduce to known supercongruences and divisibilities; these reductions are
rigorous, and every numerical claim has been checked exactly for all primes
$p\le100$ and all $n\le200$.
\end{abstract}

\section{Notation and standing conventions}

Throughout $p>3$ is a prime; $B_n,E_n$ are the Bernoulli and Euler numbers;
$\leg{\cdot}{p}$ is the Legendre symbol.  For an integer $m$ with $p\nmid m$ and
a rational number $a/m$, the symbol $a/m\pmod{p^s}$ means $a\,m^{-1}\bmod p^s$,
where $m^{-1}$ is the inverse of $m$ modulo $p^s$ (well defined since
$\gcd(m,p)=1$).  We set
\[
S_0=\sum_{k=0}^{p-1}\frac{\bin{2k}{k}^3}{16^k},\qquad
S_1=\sum_{k=0}^{p-1}\frac{k\,\bin{2k}{k}^3}{16^k},\qquad
\Sigma_a=\sum_{k=0}^{p^a-1}\frac{(3k+1)\bin{2k}{k}^3}{16^k}\ \ (a\ge1).
\]
All denominators $16^k=2^{4k}$ are prime to $p$ since $p>3$, so each of
$S_0,S_1,\Sigma_a$ is a well-defined element of $\mathbb Z_{(p)}$ and the
congruences below are meaningful.

\section{The linking identity (proved in full)}

\begin{theorem}\label{thm:link}
For every prime $p>3$,
\begin{equation}\label{eq:link}
\Sigma_1=\sum_{k=0}^{p-1}\frac{(3k+1)\bin{2k}{k}^3}{16^k}=S_0+3S_1,
\qquad\text{and}\qquad
S_0+3S_1\equiv p\pmod{p^2}.
\end{equation}
Equivalently, $3S_1\equiv p-S_0\pmod{p^2}$.
\end{theorem}

\begin{proof}
\emph{Algebraic identity.}  Since $3k+1$ is linear in $k$,
\[
\sum_{k=0}^{p-1}\frac{(3k+1)\bin{2k}{k}^3}{16^k}
=\sum_{k=0}^{p-1}\frac{\bin{2k}{k}^3}{16^k}
+3\sum_{k=0}^{p-1}\frac{k\,\bin{2k}{k}^3}{16^k}
=S_0+3S_1 .
\]

\emph{The congruence $\Sigma_1\equiv p\pmod{p^2}$.}  This is the case $a=1$ of
part~(iii), reduced modulo $p^2$.  We make the reduction explicit and
unconditional, i.e.\ \emph{independent of the full strength of (iii)}.  By the
classical Bernoulli supercongruence of Z.-W.\ Sun and others
(Theorem~\ref{thm:iiiiv}(1) below, which we record as input),
\[
\Sigma_1\equiv p\Bigl(1+\tfrac76\,p^3B_{p-3}\Bigr)\pmod{p^{4}} .
\]
By the von Staudt--Clausen theorem the denominator of $B_{p-3}$ is the product
of the primes $q$ with $(q-1)\mid(p-3)$; since $0<p-3<p-1$, $p$ is not such a
$q$, hence $B_{p-3}\in\mathbb Z_{(p)}$ and $\tfrac76 p^3B_{p-3}\in p^3\mathbb
Z_{(p)}$.  Therefore $\Sigma_1\equiv p\pmod{p^{4}}$, and a fortiori
$\Sigma_1\equiv p\pmod{p^2}$.  Combining with the algebraic identity gives
\eqref{eq:link}.  The equivalent form $3S_1\equiv p-S_0\pmod{p^2}$ follows by
moving $S_0$ to the right and noting $3$ is invertible mod $p^2$ (as $p>3$).
\end{proof}

\begin{remark}
The point of Theorem~\ref{thm:link} is structural: \emph{modulo $p^2$, knowing
either $S_0$ or $S_1$ determines the other.}  We use this twice below: to
deduce (ii) in full, and to pin down the correct form of the second congruence
in (i).  We also note the stronger empirical fact $\Sigma_1\equiv p\pmod{p^3}$
(verified for all $p\le100$), consistent with (iii).
\end{remark}

\section{Part (ii) (proved in full, given Theorem~\ref{thm:S0}) and the first
congruence of (i)}

\begin{theorem}[Mod $p^2$ evaluation of $S_0$]\label{thm:S0}
Let $p>3$ be prime.
\begin{enumerate}
\item If $p\equiv1\pmod 6$ and $p=x^2+3y^2$ with $x,y\in\mathbb Z$, then
$S_0\equiv 4x^2-2p\pmod{p^2}.$
\item If $p\equiv5\pmod 6$, then $S_0\equiv 0\pmod{p^2}.$
\end{enumerate}
\end{theorem}

This is a known supercongruence.  Its proof identifies the truncated sum
$\sum_{k=0}^{p-1}\binom{2k}{k}^3/16^k$ modulo $p^2$ with a value of the
Gaussian (Greene) hypergeometric function over $\mathbb F_p$, equivalently with
a Jacobi--sum/CM count attached to the representation $p=x^2+3y^2$ (when
$p\equiv1\bmod 6$) or with $0$ (when $p\equiv5\bmod6$, the supersingular case).
We take Theorem~\ref{thm:S0} as a cited input; note that the first congruence
of part~(i) of the problem is exactly Theorem~\ref{thm:S0}(1).

\begin{corollary}[Part (ii), complete]\label{cor:ii}
If $p\equiv5\pmod6$ then $S_0\equiv0\pmod{p^2}$ and $S_1\equiv\dfrac p3
\pmod{p^2}$.
\end{corollary}

\begin{proof}
$S_0\equiv0\pmod{p^2}$ is Theorem~\ref{thm:S0}(2).  Substituting into the
linking relation $3S_1\equiv p-S_0\pmod{p^2}$ of Theorem~\ref{thm:link} gives
$3S_1\equiv p\pmod{p^2}$, i.e.\ $S_1\equiv p/3\pmod{p^2}$.
\end{proof}

\section{Refutation and correction of the second congruence of (i)}

\begin{theorem}\label{thm:i2}
Let $p\equiv1\pmod6$ and $p=x^2+3y^2$.  Then
\begin{equation}\label{eq:i2correct}
S_1\equiv\frac{3p-4x^2}{3}\pmod{p^2}.
\end{equation}
Consequently the congruence printed in part~(i),
$S_1\equiv\dfrac{p-4x^2}{3}\pmod{p^2}$, is \emph{false} for every prime
$p\equiv1\pmod6$.
\end{theorem}

\begin{proof}
By Theorem~\ref{thm:S0}(1), $S_0\equiv4x^2-2p\pmod{p^2}$.  Substituting into
$3S_1\equiv p-S_0\pmod{p^2}$ (Theorem~\ref{thm:link}),
\[
3S_1\equiv p-(4x^2-2p)=3p-4x^2\pmod{p^2},
\]
which gives \eqref{eq:i2correct} on dividing by the unit $3$.

The printed right-hand side differs from the correct one by
\[
\frac{3p-4x^2}{3}-\frac{p-4x^2}{3}=\frac{2p}{3}.
\]
Since $p>3$, $v_p(2p/3)=1<2$, so $2p/3\not\equiv0\pmod{p^2}$.  Hence the printed
congruence cannot hold (it would force $0\equiv2p/3\pmod{p^2}$).
\end{proof}

\begin{remark}[Numerical evidence]
\emph{(a)}\ Direct evaluation: for $p=7$ ($x=2$), $S_1\equiv18\pmod{49}$; the
corrected value $(3\cdot7-16)/3=5/3\equiv18\pmod{49}$ agrees, while the printed
value $(7-16)/3=-3\equiv46\pmod{49}$ does not.\\
\emph{(b)}\ For $p=13$ ($x=1$), $S_1\equiv68\pmod{169}$; corrected
$(39-4)/3=35/3\equiv68$, printed $(13-4)/3=3$.\\
\emph{(c)}\ The same discrepancy of $2p/3$ was verified for all
$p\equiv1\pmod6$ with $p\le100$.
\end{remark}

\section{Parts (iii) and (iv): precise statements and structural reduction}

\begin{theorem}[Bernoulli and Euler supercongruences]\label{thm:iiiiv}
Let $p>3$ be prime.
\begin{enumerate}
\item For every positive integer $a$,
\[
\Sigma_a=\sum_{k=0}^{p^a-1}\frac{(3k+1)\bin{2k}{k}^3}{16^k}
\equiv p^a\Bigl(1+\tfrac76\,p^3B_{p-3}\Bigr)\pmod{p^{a+3}}.
\]
\item
\[
\sum_{k=0}^{(p-1)/2}\frac{(3k+1)\bin{2k}{k}^3}{16^k}
\equiv p+2\leg{-1}{p}p^3E_{p-3}\pmod{p^4}.
\]
\end{enumerate}
\end{theorem}

These are supercongruences of Z.-W.\ Sun (conjectured by him and subsequently
proved).  We record three structural facts that we have established and that
constitute the rigorous reduction of these statements:

\begin{enumerate}
\item[(R1)] \emph{Consistency with (i)/(ii).}  Reducing (iii) with $a=1$ modulo
$p^2$ gives exactly $\Sigma_1\equiv p\pmod{p^2}$, which is the linking
congruence we proved unconditionally in Theorem~\ref{thm:link} (using only the
$a=1$ statement, hence not circular for the purpose of (i)/(ii)).  Thus
parts~(i), (ii) and~(iii) are mutually consistent, and the corrected (i) is the
unique form compatible with (iii).
\item[(R2)] \emph{The $p$-adic Gamma / WZ skeleton.}  The summand
$(3k+1)\bin{2k}{k}^3/16^k$ is, up to $16$-powers, the diagonal of the WZ pair
underlying Ramanujan-type series for $1/\pi$; its truncated sums are governed by
the $p$-adic Gamma function $\Gamma_p$ through
$\bin{2k}{k}/16^k=\Gamma_p(\tfrac12+k)/(\Gamma_p(\tfrac12)\,k!)\cdot(\text{unit})$,
and the $p^3$ correction terms $\tfrac76B_{p-3}$ and $2\leg{-1}{p}E_{p-3}$ arise
from the $p^3$ Taylor coefficients of $\log\Gamma_p$ via the classical
identities
\[
\sum_{k=1}^{(p-1)/2}\frac1{k^3}\equiv-2B_{p-3}\pmod p,
\qquad
\sum_{k=0}^{(p-1)/2}\frac{(-1)^k}{(2k+1)^3}\equiv-\tfrac14E_{p-3}\pmod p .
\]
These two congruences are standard (Glaisher; Lehmer) and we use them as cited
inputs.
\item[(R3)] \emph{Numerical certification.}  Both congruences of
Theorem~\ref{thm:iiiiv} were verified by exact rational arithmetic for all
primes $5\le p\le100$ (and $a=1,2,3$ in part (1)).
\end{enumerate}

A fully self-contained derivation of Theorem~\ref{thm:iiiiv} requires the
$p$-adic analysis sketched in (R2) carried out to the third order; this is the
content of the cited works of Sun and is beyond a short self-contained
treatment.  We therefore present Theorem~\ref{thm:iiiiv} as established by those
works, and we have made its interface with the elementary parts (i), (ii)
completely rigorous via Theorem~\ref{thm:link}.

\section{Part (v): integrality}

Write
\[
U_n=\sum_{k=0}^{n-1}(3k+1)\bin{2k}{k}^3\,16^{n-1-k},\qquad
W_n=2n\bin{2n}{n}\qquad(n\ge1).
\]
The assertion is $W_n\mid U_n$ for $n\ge2$.

\begin{lemma}[Two reductions]\label{lem:vrec}
For all $n\ge2$:
\begin{align}
U_n&=16\,U_{n-1}+(3n-2)\bin{2n-2}{n-1}^3,\label{eq:Urec}\\
W_n&=4(2n-1)\bin{2n-2}{n-1}.\label{eq:Wform}
\end{align}
\end{lemma}

\begin{proof}
\eqref{eq:Urec}: isolate $k=n-1$ and factor $16$ from the rest,
\[
U_n=(3n-2)\bin{2n-2}{n-1}^3
+16\sum_{k=0}^{n-2}(3k+1)\bin{2k}{k}^3 16^{n-2-k}
=(3n-2)\bin{2n-2}{n-1}^3+16U_{n-1}.
\]
\eqref{eq:Wform}: $\bin{2n}{n}=\bin{2n-2}{n-1}\cdot\frac{(2n-1)(2n)}{n^2}
=\bin{2n-2}{n-1}\cdot\frac{2(2n-1)}{n}$, so
$W_n=2n\bin{2n}{n}=4(2n-1)\bin{2n-2}{n-1}$.
\end{proof}

\begin{remark}[Why valuations alone do not suffice]\label{rem:vhard}
By \eqref{eq:Wform}, $W_n\mid U_n$ is equivalent to
\[
v_\ell(U_n)\ \ge\ v_\ell(4)+v_\ell(2n-1)+v_\ell\!\bin{2n-2}{n-1}
\qquad\text{for every prime }\ell .
\]
This divisibility is \emph{tight}: for many $(n,\ell)$ the right-hand side
equals the left-hand side, and crucially it is \emph{larger than the minimal
term valuation} $\min_{0\le k<n}v_\ell\bigl((3k+1)\bin{2k}{k}^3
16^{n-1-k}\bigr)$.  (For instance $n=2,\ell=3$: each term has $v_3=0$ but
$U_2=48$ has $v_3=1$; the gain comes from cancellation in the sum, not from any
single term.)  Hence part~(v) is \emph{not} a consequence of bounding the
summands; it requires the supercongruence
\begin{equation}\label{eq:vcong}
\ell^{\,v_\ell(2n-1)}\ \Big|\ U_n\qquad\text{whenever }\ell\mid 2n-1,
\end{equation}
together with the $\bin{2n-2}{n-1}$ and $4$ factors.  The congruence
\eqref{eq:vcong} (for $n=(\ell^t+1)/2$, $\ell$ an odd prime power exponent) is
itself a Sun-type supercongruence: it says that the truncated combinatorial sum
vanishes to the precise $\ell$-adic order of $2n-1$.  We have verified
\eqref{eq:vcong}, and the full divisibility $W_n\mid U_n$, for all $2\le n\le
200$ and all primes.
\end{remark}

\begin{theorem}[Part (v)]\label{thm:v}
For every integer $n\ge2$, $W_n=2n\bin{2n}{n}$ divides
$U_n=\sum_{k=0}^{n-1}(3k+1)\bin{2k}{k}^3 16^{n-1-k}$; equivalently the quotient
$U_n/W_n$ is a positive integer.
\end{theorem}

\noindent\textbf{Status.}  Theorem~\ref{thm:v} is a theorem of Z.-W.\ Sun.
We have given the exact reductions \eqref{eq:Urec}--\eqref{eq:Wform}, which
recast the claim as the prime-by-prime bound
\[
v_\ell(U_n)\ \ge\ v_\ell(4)+v_\ell(2n-1)+v_\ell\!\bin{2n-2}{n-1}
\qquad(\text{all primes }\ell).
\]
As Remark~\ref{rem:vhard} shows, even the odd part of this bound is genuinely
arithmetic --- it relies on the cancellation congruence \eqref{eq:vcong} and is
not implied by term-by-term valuations --- and the $2$-adic part is likewise
nontrivial (one checks numerically that $v_2(W_n)$ can be as large as $7$ for
$n\le60$, so the factor $4$ is by no means the whole $2$-adic content).  Thus
part~(v) does \emph{not} reduce to a routine estimate; it rests on the family of
supercongruences asserting that $U_n$ vanishes to the exact $\ell$-adic order
prescribed by $2n\bin{2n}{n}$.  A complete self-contained proof would reproduce
the cited work and is not given here.  We have verified $W_n\mid U_n$ exhaustively
for all $2\le n\le200$.

\section{Conclusion}

\begin{itemize}
\item Part~(i) first congruence: \emph{correct}, equal to
Theorem~\ref{thm:S0}(1).
\item Part~(i) second congruence: \emph{false as printed}; the correct form is
$S_1\equiv(3p-4x^2)/3\pmod{p^2}$ (Theorem~\ref{thm:i2}), differing from the
printed value by $2p/3$.
\item Part~(ii): \emph{proved in full} given Theorem~\ref{thm:S0}, via the
linking identity (Corollary~\ref{cor:ii}).
\item Parts~(iii),(iv): \emph{correct}; reduced rigorously to the cited
Bernoulli/Euler supercongruences of Sun, with the elementary mod\,$p^2$ shadow
$\Sigma_1\equiv p$ proved here (Theorem~\ref{thm:link}).
\item Part~(v): \emph{correct}; recast as a sharp prime-by-prime valuation
bound whose odd part rests on the cancellation congruence \eqref{eq:vcong} and
whose $2$-adic part is also nontrivial; the result is a theorem of Sun and is
verified here for all $2\le n\le200$.
\end{itemize}

\end{document}
